% Direct-authority diplomatic French for printed p.208.
\clearpage
\SourcePage{208}

\LayerTarget{tohoku:proof:lemma-5.5.2:p1}
\noindent ctérisation des objets injectifs dans $\mathbf K(C_1)$ indiquée dans
2.4, $C$ est acyclique en dimensions $>0$ et se ``décompose'', de plus les
$C^i$ et $Z^i(C)$ sont injectifs, d'où il résulte d'abord que $\overline F(C)$
est une résolution de $F(Z^0(C))$ d'où
$R^nG(\overline F(C))=R^nG(F(Z^0(C)))$, et comme $Z^0(C)$ est injectif le
deuxième membre est nul d'après l'hypothèse sur $F$. Ainsi, on peut appliquer
le théorème 2.4 au foncteur composé $(H^0\overline G)\overline F$, ce qui donne
le foncteur spectral de l'énoncé, compte tenu que
$R^i\overline F=\overline{R^iF}$ (relation qui résulte immédiatement des
définitions).

\LayerTarget{tohoku:corollary:lemma-5.5.2}
\noindent\textsc{Corollaire.} \textit{Sous les conditions précédentes, soit $C$
une résolution d'un objet $A\in C_1$. Alors il existe une suite spectrale
cohomologique, aboutissant à l'objet gr dué $(R^n(GF)(A))$, et dont le terme
initial est donné dans (5.5.2).}

\LayerTarget{tohoku:proof:corollary-lemma-5.5.2}
En effet, on a ici $R^n(GF)(C)=R^n(GF)(A)$.

\LayerTarget{tohoku:definition:g-cover}
Revenons à l'espace $X$ avec le groupe d'opérateurs $G$. Soit
$\mathcal U=(U_i)_{i\in I}$ un recouvrement de $X$, $G$ opérant dans l'ensemble
d'indices $I$ de telle façon qu'on ait $g\cdot U_i=U_{g\cdot i}$ pour tout
$i\in I$ : on dira que $\mathcal U$ est un $G$-recouvrement. Si $P$ est un
préfaisceau abélien sur $X$ dans lequel $G$ opère, alors le complexe
$C(\mathcal U,P)$ des cochaînes de $\mathcal U$ à coefficients dans $P$ peut
être considéré comme un complexe de $G$-groupes de façon évidente. Nous
poserons

\LayerTarget{tohoku:equation:5.5.3}
\[
\tag{5.5.3}
H^n(\mathcal U;G,P)=R^n\Gamma^G(C(\mathcal U,P))
\]

\LayerTarget{tohoku:construction:equivariant-cech-bicomplex}
où le deuxième membre peut donc aussi s'expliciter comme le groupe
$H^n(C(G,C(\mathcal U,P)))$, $C(G,C(\mathcal U,P))$ étant le bicomplexe formé
des cochaînes de $G$ à valeurs dans le complexe $C(\mathcal U,P)$, bicomplexe
qu'on pourra noter $C(\mathcal U;G,P)$ :

\LayerTarget{tohoku:equation:5.5.3-bis}
\[
\tag{5.5.3 bis}
H^n(\mathcal U;G,P)=H^n(C(\mathcal U;G,P)).
\]

\LayerTarget{tohoku:construction:open-g-cover-spectral-functor}
Supposons d'abord le $G$-recouvrement $\mathcal U$ ouvert. Nous allons appliquer
le lemme 5.5.1 en prenant les catégories $C^{X(G)}$, $C^G$ (catégorie des
$G$-modules), $C$ (catégorie des groupes abéliens), et les foncteurs
$\Gamma_X$ et $\Gamma^G$, et posant enfin $\mathbf K(A)=C(\mathcal U,A)$.
$\Gamma_X$ transforme bien objets injectifs en objets $\Gamma^G$-acycliques
(corollaire à la prop. 5.1.3), et les conditions (i) et (ii) du lemme sont
satisfaites, (i) parce que $\mathcal U$ est ouvert, et (ii) en vertu de ce qui a
été dit dans 3.8 (deuxième alinéa). Il reste à expliciter $R^q\mathbf K(A)$, le
calcul est immédiat (et a d'ailleurs été fait dans 3.8) on trouve
$R^q\mathbf K(A)=C(\mathcal U,H^q(A))$, où $H^q(A)$ désigne le préfaisceau
$H^q(A)(V)=H^q(V,A)$. On a donc prouvé :

\LayerTarget{tohoku:theorem:5.5.3}
\noindent\textsc{Théorème 5.5.3.} \textit{Soit $X$ un espace muni d'un groupe
$G$ d'opérateurs, soit $\mathcal U=(U_i)$ un $G$-recouvrement ouvert de $X$.
Alors il existe sur la catégorie $C^{X(G)}$ des $G$-faisceaux abéliens sur $X$,
un foncteur spectral cohomologique, aboutissant au foncteur gradué
$(H^n(X;G,A))$, et dont le terme initial est}

\LayerTarget{tohoku:equation:5.5.4}
\[
\tag{5.5.4}
E_2^{p,q}(A)=H^p(\mathcal U;G,H^q(A))
\]
\textit{(où $H^q(A)$ est le préfaisceau $H^q(A)(V)=H^q(V,A)$ sur $X$).}

\LayerTarget{tohoku:map:5.5.5}
On en conclut comme d'habitude des ``edge-homomorphisms''
\[
\tag{5.5.5}
H^n(\mathcal U;G,A)\longrightarrow H^n(X;G,A)
\]

\LayerTarget{tohoku:conclusion:theorem-5.5.3:p1}
et une suite exacte à 5 termes, montrant entre autres que l'homomorphisme

\TohokuContinuation{209}
